\documentclass[11pt]{article}
\usepackage[a4paper,margin=25mm]{geometry}
\usepackage{amsmath,amssymb,amsthm}
\newtheorem{theorem}{Theorem}
\newtheorem{proposition}[theorem]{Proposition}
\title{Chiral quotients from one parity subgroup}
\author{}\date{}
\begin{document}\maketitle
We recover the quotient criterion and search of Hartley, Hubard and
Leemans, \emph{Two atlases of abstract chiral polytopes for small groups},
Ars Math. Contemp. 5 (2012), printed page 375. The underlying group-action
argument is classical. Its role here is to specify one reusable interface,
including the two different equivalence conventions in that paper.
This note does not establish either of the paper's complete finite atlases.

Let $W=\langle r_0,\ldots,r_{n-1}\rangle$ be a finite group generated by
involutions, and $N\leq W$. On the right cosets $Nw$, put a colour-$i$
edge from $Nw$ to $Nwr_i$. Loops and coincident edges are permitted at
this algebraic stage. The graph is connected because the $r_i$ generate.
When a quotient polytope is used below, we require its actual flag graph
to be this graph. The source's semisparse quotient theorem supplies that
geometric interface; an arbitrary orbit poset is not sufficient.

\begin{theorem}\label{normalizer}
The colour-preserving automorphism group of the coset graph is
$N_W(N)/N$. It acts freely, with $[W:N_W(N)]$ vertex orbits.
If this is the flag graph of a polytope, that polytope is chiral if and
only if $[W:N_W(N)]=2$ and $r_i\notin N_W(N)$ for every $i$.
\end{theorem}
\begin{proof}
A colour-preserving automorphism commutes with each of the permutations
induced by $r_i$, hence with their products. If it sends $N$ to $Na$,
then it must send $Nw$ to $Naw$. For $x\in N$, well-definedness at $N$
requires $Nax=Na$, so $aNa^{-1}\leq N$. The two subgroups have the same
finite order, hence equality. Thus $a\in N_W(N)$. Conversely any such
$a$ gives a well-defined bijection $Nw\mapsto Naw$ commuting with all
colours. Two choices give the same map exactly when their cosets modulo
$N$ agree. This proves the asserted group identification.

If $Naw=Nw$, then $a\in N$, so the action is free. Its orbits therefore
number
\[
 (|W|/|N|)/(|N_W(N)|/|N|)=[W:N_W(N)].
\]
A polytope's automorphisms
are precisely its colour-preserving flag-graph automorphisms: its
$i$-faces are the components after deleting colour $i$, and face incidence
is witnessed by flags. This uses the stipulated actual flag graph.

Write $H=N_W(N)$ and suppose $[W:H]=2$, so $H\lhd W$. The vertices
$Nw$ and $Nwr_i$ are in the same $H/N$-orbit exactly when
$wr_iw^{-1}\in H$, equivalently $r_i\in H$. Consequently every adjacency
exchanges the two orbits exactly when all $r_i$ lie outside $H$.
This is the defining flag-orbit condition for chirality.
\end{proof}

\begin{proposition}\label{search}
There is at most one subgroup $H$ of index two avoiding every $r_i$.
When it exists, the chiral quotients in Theorem~\ref{normalizer} arise
precisely from the subgroups $N\lhd H$ that are not normal in $W$ and
meet the geometric interface. Their conjugacy classes in $W$ have two
members, $N$ and $r_0Nr_0$. The kernel of the flag action of $W$ is
\[
 C=\operatorname{Core}_W(N)=N\cap r_0Nr_0.
\]
In $H/C$, the images of these two normal subgroups commute and intersect
trivially. In particular, if the flag action is faithful their product
is an internal direct product.
\end{proposition}
\begin{proof}
Every such index-two subgroup is the kernel of the homomorphism to
$C_2$ sending each $r_i$ to its nonidentity element. A map on a generating
set determines a homomorphism uniquely. This also gives an effective
existence test: assign parity to vertices of the Cayley graph, and reject
if any edge contradicts the assignment. Equivalently, the subgroup
generated by the $r_0r_i$ has index two; otherwise no such $H$ exists.
Indeed every even word is a product of pairs, and
$r_ir_j=(r_0r_i)^{-1}(r_0r_j)$.

If $N\lhd H$ is not normal in $W$, its normalizer contains $H$ but is
not $W$, and therefore equals $H$. Conversely the normalizer condition
of Theorem~\ref{normalizer} implies exactly these requirements. Since
$W=H\sqcup Hr_0$, conjugation produces just the two displayed subgroups.
The usual transitive coset-action calculation gives its kernel as the
intersection of all conjugates, hence the displayed $C$.

Both $N$ and $r_0Nr_0$ are normal in $H$. The commutator of an element
of one with an element of the other belongs to both, hence to $C$.
Their images in $H/C$ therefore commute and have trivial intersection;
the product map from their direct product is an isomorphism onto its
image. Finally, coset graphs of subgroups $N,M\leq W$ are colour-isomorphic
exactly when $N,M$ are conjugate: the same argument as above, for a map
$N\mapsto Ma$, gives $aNa^{-1}=M$. Thus the two conjugates describe one
unrooted quotient, rather than two different polytopes.
\end{proof}

The condition $C=1$ is exactly faithfulness of this flag action. We do
not infer that an arbitrary monodromy quotient is itself a regular
polytope, or identify all notions of minimal polytopal cover. The source's
qualified cover statement must be read with this distinction. Its
small-cover algorithm removes the conjugate pair just proved. The
almost-simple atlas instead reports both enantiomorphic markings before
dividing by two; the acquired prefix-coordinate theorem treats those
as oriented and unrooted equivalence, respectively.

\paragraph{A literal finite instance.}
The following instance exercises the interface without assuming any
target atlas count. Over $\mathbf F_7$, put $V=\mathbf F_7^2$,
$R=\operatorname{diag}(3,5)$, and let $S$ exchange the two coordinates.
Then $R$ has order six and $SRS=R^{-1}$. In
$W=V\rtimes\langle R,S\rangle$, of order $588$, use
\[
 r_0=((1,1),R^3S),\qquad r_1=(0,RS),\qquad r_2=(0,S).
\]
These are involutions with pair orders $3,6,2$, in string order.
They generate: $r_1r_2=R$, while $r_0(R^3S)^{-1}$ is translation by
$(1,1)$; this vector and its image $(3,5)$ under $R$ span $V$.
The parity subgroup is $H=V\rtimes\langle R\rangle$.

The only normal subgroups of $H$ that are not normal in $W$ are the
two coordinate lines. To see this, if a normal subgroup contains an
element with rotation part $R^k\ne1$, its commutators with $V$ contain
$(R^k-I)V=V$, because both diagonal entries of $R^k-I$ are nonzero.
It is then the full inverse image of a subgroup of $H/V=C_6$, and is
normal in $W$, since $S$ acts on that cyclic quotient by inversion.
Otherwise it is an $R$-invariant additive subgroup of $V$. Such a
subgroup is an $\mathbf F_7$-subspace. The projections
$(R-5I)/(3-5)$ and $(R-3I)/(5-3)$ show that it is a sum of coordinate
lines. Of the four possibilities, only the two lines are not $S$-invariant.

The literal checker reconstructs all $588$ group elements, verifies the
full intersection property of the displayed regular marking, and builds
the complete $84$-flag quotient for either line. It checks the entire
face poset, its flags, diamond intervals and every flag residue, not
just the orders of adjacent products. Independently propagating every
possible image of one vertex through the coloured graph gives exactly
$42$ automorphisms, agreeing with $H/N$. The lines are exchanged by
$r_0$, intersect trivially, and yield one unrooted chiral quotient of
this particular regular cover. Its mirror pair is not counted twice.
This is a bounded validation example, not a proof of the source's
complete $48+8$ or $17114$ enumerations.
\end{document}
